\chapter{نتیجه‌گیری}
\label{ch:conclusion}

این پایان‌نامه یک ارزیابی تمام‌کوهورت و گذشته‌نگر از پیش‌بینی نشانگر جانشین امکان‌پذیری ترخیص از \lr{ICU} در \lr{MIMIC\text{-}IV v3.1} ارائه کرد. نسخۀ نهایی از \pnum{32,399} بستری نخست بزرگسال استفاده کرد و دیگر بر زیرمجموعۀ \pnum{2,000} بستری به‌عنوان تحلیل اصلی متکی نیست.

پرسش نخست پژوهش پاسخ مثبت مشروط گرفت: داده‌های \pnum{24} ساعت نخست \lr{ICU} توان پیش‌بینی متوسط تا خوب برای برچسب جانشین داشتند. پرسش دوم نشان داد مدل‌های جدولی و درختی از \lr{Few-Shot ETHOS} خام قوی‌تر بودند. پرسش سوم نشان داد توکن‌های علائم در حالت هیبریدی مفیدتر از حالت مستقل‌اند. پرسش چهارم نشان داد فدرال شبیه‌سازی‌شده می‌تواند به خط پایۀ لجستیک متمرکز نزدیک بماند. پرسش پنجم نشان داد \lr{XGBoost} بهترین توازن عملی میان تفکیک و کالیبراسیون را فراهم کرد، هرچند مدل انباشتی بالاترین \lr{AUROC} را داشت.

جمع‌بندی علمی کار این است که نوآوری روش‌شناختی باید در برابر خط پایه‌های قوی آزموده شود. در این وظیفۀ ساخت‌یافتۀ \lr{EHR}، روش چندنمونه‌ای خام برتری نداشت؛ اما همین نتیجه پایان‌نامه را صادقانه‌تر و مفیدتر می‌کند. مدل‌های درختی، به‌ویژه \lr{XGBoost}، برای داده‌های جدولی بالینی بسیار رقابتی باقی ماندند. توکن‌های علائم و یادگیری فدرال نیز مسیرهای پژوهشی قابل‌پیگیری هستند، به شرط آنکه در مطالعات آینده با برچسب‌های بالینی دقیق‌تر، داده‌های چندمرکزی و اعتبارسنجی آینده‌نگر ارزیابی شوند.

بنابراین، دستاورد اصلی پایان‌نامه یک سامانه آماده استقرار نیست؛ بلکه یک بنچمارک تمام‌کوهورت، شفاف و بازتولیدپذیر است که نشان می‌دهد در مسئلۀ پیش‌بینی جانشین امکان‌پذیری ترخیص از \lr{ICU}، کدام رویکردها اکنون قوی‌ترند و کدام مسیرها برای پژوهش‌های بعدی باید بهبود یابند.
